% ============================================================
\chapter{Compacit\'e dans les Espaces Non-Hausdorff}
\label{ch:compacite}
% ============================================================

La compacit\'e est l'une des notions les plus rassurantes de la topologie classique~: dans un espace de Hausdorff, un compact est ferm\'e, born\'e (en m\'etrique), et poss\`ede toutes les propri\'et\'es de finitude que l'on peut souhaiter. Mais que se passe-t-il quand on abandonne l'axiome de Hausdorff~? Les compacts cessent d'\^etre ferm\'es, les points limites ne sont plus uniques, et l'intuition g\'eom\'etrique vacille. Ce n'est pas un accident~: en th\'eorie des domaines, les espaces qui mod\'elisent les calculs informatiques sont typiquement $T_0$ mais pas $T_1$. La notion de \emph{compact satur\'e} --- un compact \'egal \`a l'intersection de ses voisinages ouverts --- remplace alors celle de compact ferm\'e, et la compactification de Smyth joue le r\^ole que l'Alexandroff jouait en Hausdorff. Ce chapitre d\'eveloppe cette th\'eorie raffin\'ee.

\begin{intuition}
Dans les espaces de Hausdorff, la compacit\'e est un concept bien compris~: les compacts
sont ferm\'es, l'intersection de compacts est compacte, et la compactification d'Alexandroff
produit un espace de Hausdorff si et seulement si l'espace de d\'epart est localement compact
et s\'epar\'e. En revanche, dans les espaces non-Hausdorff, la compacit\'e pr\'esente des
subtilit\'es fondamentales. Les compacts ne sont plus n\'ecessairement ferm\'es, et il faut
recourir \`a la notion de \emph{compact satur\'e} pour retrouver un comportement raisonnable.
Ce chapitre d\'eveloppe la th\'eorie de la compacit\'e adapt\'ee aux espaces $T_0$, avec
des applications en th\'eorie des domaines.
\end{intuition}

% ------------------------------------------------------------
\section{Compacts satur\'es}
% ------------------------------------------------------------

\begin{definition}[Ensemble satur\'e]
\index{ensemble satur\'e}
Soit $(X,\tau)$ un espace topologique. Un sous-ensemble $A \subseteq X$ est \emph{satur\'e}
s'il est \'egal \`a l'intersection de tous les ouverts qui le contiennent~:
\[
  A = \bigcap \{ U \in \tau : A \subseteq U \}.
\]
De mani\`ere \'equivalente, $A$ est satur\'e si et seulement si $A$ est une partie haute pour
l'ordre de sp\'ecialisation~: si $x \in A$ et $x \leq y$ (c'est-\`a-dire $x \in \overline{\{y\}}$),
alors $y \in A$.
\end{definition}

\begin{remarque}
Dans un espace $T_1$, tout sous-ensemble est satur\'e, car l'ordre de sp\'ecialisation est
l'\'egalit\'e. La notion de saturation n'est donc pertinente que dans les espaces non $T_1$.
\end{remarque}

\begin{proposition}[Saturation d'un ensemble]
Pour tout sous-ensemble $A$ d'un espace topologique $X$, la \emph{saturation} de $A$ est
\[
  \uparrow\! A = \{ y \in X : \exists x \in A, \; x \leq y \}
\]
o\`u $\leq$ d\'esigne l'ordre de sp\'ecialisation. C'est le plus petit ensemble satur\'e
contenant $A$.
\end{proposition}

\begin{definition}[Compact satur\'e]
\index{compact satur\'e}
Un sous-ensemble $Q$ d'un espace topologique est un \emph{compact satur\'e} s'il est
\`a la fois compact et satur\'e.
\end{definition}

\begin{attention}
Dans un espace non-Hausdorff, un compact n'est pas n\'ecessairement ferm\'e. Par exemple,
dans l'espace de Sierpi\'nski $\mathbb{S} = \{0,1\}$ avec $\tau = \{\emptyset, \{1\}, \{0,1\}\}$,
le singleton $\{1\}$ est compact (fini) et ouvert, mais n'est pas ferm\'e. En revanche,
il est satur\'e.
\end{attention}

\begin{theoreme}[Compacts satur\'es et intersection]
\label{thm:inter-compact-sature}
Soit $X$ un espace topologique. L'intersection d'un compact satur\'e $Q$ et d'un ferm\'e
compact $K$ est compacte. Plus g\'en\'eralement, l'intersection finie de compacts satur\'es
est un compact satur\'e.
\end{theoreme}

\begin{proof}
Soit $(U_i)_{i \in I}$ un recouvrement ouvert de $Q_1 \cap Q_2$, o\`u $Q_1$ et $Q_2$
sont deux compacts satur\'es. Pour tout $x \in Q_1 \setminus Q_2$, comme $Q_2$ est satur\'e,
il existe un ouvert $V_x$ contenant $Q_2$ mais pas $x$. Alors
$(U_i)_{i \in I} \cup (X \setminus Q_2)$ recouvre $Q_1$, et par compacit\'e de $Q_1$, on
extrait un sous-recouvrement fini. Les ouverts de ce sous-recouvrement qui sont parmi les
$U_i$ recouvrent $Q_1 \cap Q_2$.
\end{proof}

\begin{exemple}
Consid\'erons l'ensemble $\N$ muni de la topologie d'Alexandrov associ\'ee \`a l'ordre
usuel. Les ouverts sont les parties hautes $\{n, n+1, n+2, \ldots\}$ et $\emptyset$.
Tout ensemble fini $\{n\}$ est compact mais pas satur\'e (car $\uparrow\! \{n\} = \{n, n+1, \ldots\}$).
Les compacts satur\'es sont les ensembles $\{n, n+1, \ldots\}$ qui sont finis ou \'egaux
\`a $\N$ tout entier~--- mais en fait $\{n, n+1, \ldots\}$ n'est pas compact. Ainsi le seul
compact satur\'e non vide est $\N$ lui-m\^eme.
\end{exemple}

% ------------------------------------------------------------
\section{Compactification d'Alexandroff pour les espaces non-Hausdorff}
% ------------------------------------------------------------

\begin{definition}[Compactification d'Alexandroff]
\index{compactification d'Alexandroff}
Soit $(X,\tau)$ un espace topologique. La \emph{compactification par un point} (ou
d'Alexandroff) est l'espace $X^* = X \cup \{\infty\}$ muni de la topologie
\[
  \tau^* = \tau \cup \{ X^* \setminus K : K \subseteq X \text{ compact et satur\'e} \}.
\]
\end{definition}

\begin{theoreme}[Propri\'et\'es de la compactification d'Alexandroff]
Soit $X$ un espace topologique.
\begin{enumerate}
  \item $X^*$ est compact.
  \item L'inclusion $X \hookrightarrow X^*$ est un plongement topologique (ouvert).
  \item Si $X$ est $T_0$, alors $X^*$ est $T_0$.
  \item Si $X$ est sobre, alors $X^*$ est sobre si et seulement si $X$ est localement compact.
\end{enumerate}
\end{theoreme}

\begin{proof}
(1) Soit $(U_i)_{i \in I}$ un recouvrement ouvert de $X^*$. L'un des $U_i$ contient
$\infty$, donc est de la forme $X^* \setminus K$ pour un compact satur\'e $K$. Les
ouverts restants recouvrent $K$, et par compacit\'e de $K$, on en extrait un
sous-recouvrement fini.

(2) Les ouverts de $X$ dans $X^*$ sont exactement les \'el\'ements de $\tau$, car
$U \cap X = U$ pour $U \in \tau$.

(3) Si $X$ est $T_0$ et $x \in X$, alors $\{x\}$ et $\{\infty\}$ sont s\'epar\'es par
l'ouvert $X$ de $X^*$.

(4) Admis (voir la section suivante pour les espaces sobres localement compacts).
\end{proof}

\begin{remarque}
Contrairement au cas de Hausdorff, la compactification d'Alexandroff d'un espace non-Hausdorff
n'est jamais de Hausdorff (sauf si l'espace de d\'epart l'est d\'ej\`a et est localement
compact). N\'eanmoins, elle pr\'eserve la $T_0$-s\'eparation, ce qui est le niveau de
s\'eparation naturel en th\'eorie des domaines.
\end{remarque}

% ------------------------------------------------------------
\section{Espaces sobres localement compacts}
% ------------------------------------------------------------

\begin{definition}[Espace localement compact]
\index{espace localement compact}
Un espace topologique $X$ est \emph{localement compact} si pour tout $x \in X$ et tout
ouvert $U$ contenant $x$, il existe un ouvert $V$ et un compact satur\'e $Q$ tels que
$x \in V \subseteq Q \subseteq U$.
\end{definition}

\begin{proposition}
Dans un espace sobre, la compacit\'e locale est \'equivalente \`a la condition suivante~:
le treillis des ouverts $\mathcal{O}(X)$ est un treillis continu.
\end{proposition}

\begin{definition}[Relation \og{}bien en dessous\fg]
\index{relation bien en dessous}
Dans un treillis complet $L$, on dit que $a$ est \emph{bien en dessous} de $b$, not\'e
$a \ll b$, si pour toute famille dirig\'ee $D$ telle que $b \leq \bigvee D$, il existe
$d \in D$ tel que $a \leq d$.
\end{definition}

\begin{theoreme}[Caract\'erisation des espaces sobres localement compacts]
\label{thm:sobre-lc}
Pour un espace topologique $X$, les conditions suivantes sont \'equivalentes~:
\begin{enumerate}
  \item $X$ est sobre et localement compact.
  \item Le treillis $\mathcal{O}(X)$ est un treillis continu.
  \item $X$ est hom\'eomorphe au spectre d'un treillis continu.
\end{enumerate}
\end{theoreme}

\begin{proof}
$(1) \Rightarrow (2)$~: Si $X$ est localement compact, pour tout ouvert $U$ et tout
$x \in U$, il existe $V$ ouvert et $Q$ compact satur\'e avec $x \in V \subseteq Q \subseteq U$.
On v\'erifie que $V \ll U$ dans $\mathcal{O}(X)$ et donc $U = \bigvee \{V : V \ll U\}$.

$(2) \Rightarrow (3)$~: Par le th\'eor\`eme de repr\'esentation des treillis continus, le
spectre de $\mathcal{O}(X)$ est sobre et localement compact, et $\mathcal{O}(X)$ est
isomorphe \`a son treillis d'ouverts.

$(3) \Rightarrow (1)$~: Le spectre d'un treillis continu est sobre par d\'efinition et
localement compact car la relation $\ll$ fournit les voisinages compacts satur\'es requis.
\end{proof}

% ------------------------------------------------------------
\section{Espaces stablement compacts}
% ------------------------------------------------------------

\begin{definition}[Espace stablement compact]
\index{espace stablement compact}
Un espace topologique $X$ est \emph{stablement compact} s'il v\'erifie~:
\begin{enumerate}
  \item $X$ est $T_0$~;
  \item $X$ est compact~;
  \item $X$ est localement compact~;
  \item $X$ est \emph{coh\'erent}~: l'intersection de deux compacts satur\'es est un
    compact satur\'e~;
  \item $X$ est sobre.
\end{enumerate}
\end{definition}

\begin{intuition}
Les espaces stablement compacts sont les analogues non-Hausdorff des espaces compacts de
Hausdorff. Ils apparaissent naturellement comme espaces de mod\`eles en s\'emantique
d\'enotationnelle (domaines de Scott) et en logique (espaces de Stone g\'en\'eralis\'es).
\end{intuition}

\begin{exemple}[Spectre d'un treillis distributif]
Si $L$ est un treillis distributif born\'e, son spectre (ensemble des id\'eaux premiers muni
de la topologie de Hull-Kernel) est un espace stablement compact.
\end{exemple}

\begin{theoreme}[Caract\'erisation par les cadres]
Un espace $T_0$ est stablement compact si et seulement si son treillis d'ouverts
$\mathcal{O}(X)$ est un treillis continu et $X$ est compact et coh\'erent.
\end{theoreme}

\begin{proposition}[Stabilit\'e de la compacit\'e stable]
La classe des espaces stablement compacts est stable par~:
\begin{enumerate}
  \item produit fini~;
  \item sous-espace ferm\'e satur\'e~;
  \item passage \`a la topologie patch.
\end{enumerate}
\end{proposition}

% ------------------------------------------------------------
\section{La topologie patch}
% ------------------------------------------------------------

\begin{definition}[Topologie patch (ou de Lawson)]
\index{topologie patch}
\index{topologie de Lawson}
Soit $X$ un espace topologique. La \emph{topologie patch} sur $X$ est la topologie
engendr\'ee par les ouverts de $X$ et les compl\'ementaires des compacts satur\'es de $X$.
On note $X^{\mathrm{patch}}$ l'espace r\'esultant.
\end{definition}

\begin{theoreme}[Patch d'un espace stablement compact]
\label{thm:patch}
Si $X$ est stablement compact, alors~:
\begin{enumerate}
  \item $X^{\mathrm{patch}}$ est compact de Hausdorff.
  \item Les ouverts de $X$ forment exactement les ouverts de $X^{\mathrm{patch}}$
    qui sont des parties hautes pour l'ordre de sp\'ecialisation de $X$.
  \item L'ordre de sp\'ecialisation de $X$ est un ordre partiel ferm\'e dans
    $(X^{\mathrm{patch}})^2$.
\end{enumerate}
\end{theoreme}

\begin{proof}
(1) La topologie patch est plus fine que la topologie originale, donc s\'epar\'ee
(puisque les compacts satur\'es permettent de s\'eparer les points). La compacit\'e
d\'ecoule du th\'eor\`eme d'Alexander appliqu\'e \`a la sous-base form\'ee des ouverts
et des compl\'ementaires de compacts satur\'es, en utilisant la coh\'erence de $X$.

(2) et (3) d\'ecoulent de la construction de la topologie patch et des propri\'et\'es de
l'ordre de sp\'ecialisation.
\end{proof}

\begin{exemple}
L'intervalle $[0,1]$ muni de la topologie de Scott (ouverts~: $\emptyset$ et les
intervalles $(a, 1]$ pour $a \in [0,1)$, plus $[0,1]$) est stablement compact.
Sa topologie patch est la topologie usuelle sur $[0,1]$.
\end{exemple}

% ------------------------------------------------------------
\section{Dualit\'e de de Groot}
% ------------------------------------------------------------

\begin{definition}[Topologie co-compacte (ou duale de de Groot)]
\index{dualit\'e de de Groot}
Soit $X$ un espace topologique. La \emph{topologie duale de de Groot} $\tau^d$ est la
topologie engendr\'ee par les compl\'ementaires des compacts satur\'es de $(X,\tau)$~:
\[
  \tau^d = \langle \{ X \setminus Q : Q \text{ compact satur\'e de } (X,\tau) \} \rangle.
\]
On note $X^d = (X, \tau^d)$ l'espace dual.
\end{definition}

\begin{theoreme}[Dualit\'e de de Groot pour les espaces stablement compacts]
\label{thm:de-groot}
Si $X$ est stablement compact, alors~:
\begin{enumerate}
  \item $X^d$ est stablement compact.
  \item $(X^d)^d = X$ (la dualit\'e est involutive).
  \item L'ordre de sp\'ecialisation de $X^d$ est l'ordre oppos\'e de celui de $X$.
  \item La topologie patch de $X$ et celle de $X^d$ co\"incident.
\end{enumerate}
\end{theoreme}

\begin{proof}
(1) Les compacts satur\'es de $X^d$ sont exactement les ferm\'es de $X$ qui sont des
parties basses pour l'ordre de sp\'ecialisation, c'est-\`a-dire les compacts de la topologie
originale qui sont des parties basses. On v\'erifie les cinq axiomes un par un.

(2) D\'ecoule du fait que la topologie patch d\'etermine \`a la fois $\tau$ et $\tau^d$.

(3) L'ordre de sp\'ecialisation de $\tau^d$ est d\'etermin\'e par~: $x \leq^d y$ ssi pour
tout compact satur\'e $Q$ de $\tau$, $y \in Q \Rightarrow x \in Q$, ce qui correspond
\`a $y \leq x$ dans l'ordre de sp\'ecialisation de $\tau$.

(4) La topologie patch de $X^d$ est engendr\'ee par $\tau^d$ et les compl\'ementaires
des compacts satur\'es de $X^d$, lesquels sont les ouverts de $X$. Donc elle co\"incide
avec la topologie patch de $X$.
\end{proof}

\begin{formules}[title=\textbf{R\'esum\'e~: dualit\'e de de Groot}]
\begin{center}
\renewcommand{\arraystretch}{1.3}
\begin{tabular}{|l|c|c|}
\hline
& \textbf{Espace} $X$ & \textbf{Dual} $X^d$ \\
\hline
Ouverts & $\tau$ & Engendr\'es par $X \setminus Q$, $Q$ compact satur\'e \\
Ordre de sp\'ec. & $\leq$ & $\geq$ \\
Topologie patch & $\tau^{\mathrm{patch}}$ & $(\tau^d)^{\mathrm{patch}} = \tau^{\mathrm{patch}}$ \\
\hline
\end{tabular}
\end{center}
\end{formules}

% ------------------------------------------------------------
\section{Applications en th\'eorie des domaines}
% ------------------------------------------------------------

\begin{theoreme}[Domaines comme espaces stablement compacts]
Soit $D$ un domaine continu born\'e complet (c'est-\`a-dire un dcpo continu dans lequel
tout sous-ensemble born\'e admet un supremum). Alors $D$ muni de la topologie de Scott
est stablement compact.
\end{theoreme}

\begin{corollaire}
Le domaine des intervalles $\mathbf{IR} = \{[a,b] : a \leq b, \; a,b \in \R\}$ ordonn\'e
par inclusion inverse est stablement compact pour la topologie de Scott.
\end{corollaire}

\begin{exercice}
Montrer que dans un espace $T_0$, un ensemble compact est satur\'e si et seulement s'il
est une intersection d'ouverts.
\end{exercice}

\begin{exercice}
Soit $X$ un espace sobre localement compact. Montrer que pour tout ouvert $U$ de $X$,
l'ensemble $U$ est la r\'eunion des int\'erieurs de compacts satur\'es contenus dans $U$.
\end{exercice}

\begin{exercice}
Soit $P$ un ensemble ordonn\'e fini muni de la topologie d'Alexandrov. D\'eterminer les
compacts satur\'es de $P$. En d\'eduire la topologie patch et la topologie duale de de Groot.
\end{exercice}

\begin{exercice}
Montrer que si $X$ est stablement compact et $f\colon X \to Y$ est continue, surjective
et ouverte, alors $Y$ est stablement compact.
\end{exercice}

\begin{exercice}[Dualit\'e de Stone g\'en\'eralis\'ee]
Soit $L$ un treillis distributif born\'e et $X = \mathrm{Spec}(L)$ son spectre. Montrer
que $X$ est stablement compact et que $L$ est isomorphe au treillis des compacts satur\'es
ouverts de $X$.
\end{exercice}

\begin{exercice}
Soit $X$ un espace stablement compact et $Q_1, Q_2$ deux compacts satur\'es. Montrer que
$Q_1 \cup Q_2$ et $Q_1 \cap Q_2$ sont des compacts satur\'es, et que l'ensemble des
compacts satur\'es de $X$ forme un treillis distributif.
\end{exercice}
